\documentclass{article}
\usepackage[T1]{fontenc}
\usepackage{lmodern}
\usepackage{graphicx}
\usepackage{amsmath,amssymb}
\usepackage[a4paper,margin=2cm]{geometry}
\usepackage[absolute,overlay]{textpos}
\setlength{\TPHorizModule}{1mm}
\setlength{\TPVertModule}{1mm}
\usepackage[indonesian]{babel}
\usepackage{hyperref}
\setlength{\emergencystretch}{2em}
% unit: Halaman 9

\begin{document}

% === Halaman 9 (python/native) ===

9

Properti algoritma ElGamal:

1.  Bilangan prima, p 


(tidak rahasia)

2.  Bilangan acak, g  ( g < p, g adalah akar primitif dari p)  (tidak rahasia)

3. Bilangan acak, x  (2  x  p - 2) 

(rahasia, kunci privat)

4.  y = gx mod p 



(tidak rahasia, kunci publik)

5.  m   (plainteks)  


(rahasia)

6.  a dan b  (cipherteks)  

(tidak rahasia)

\end{document}
